%%
% BIThesis 本科毕业设计论文模板 —— 使用 XeLaTeX 编译 The BIThesis Template for Undergraduate Thesis
% This file has no copyright assigned and is placed in the Public Domain.
%%

% 致谢部分尽量不使用 \subsection 二级标题，只使用 \section 一级标题
\begin{acknowledgements}
  转眼间，四年的本科学习生活即将画上句号。回望入学时的青涩与此刻略增的笃定，我深知成长离不开师长、家人与同伴的扶持与包容。在此，谨向所有帮助过我的人们致以诚挚谢意。

  首先要衷心感谢导师方浩教授。方老师在选题凝练、文献脉络与技术路线取舍上耐心点拨，并在论文结构与表述细节上严格要求、反复批注，帮助我建立更条理的写作习惯；在实验受挫与文稿屡次返工的阶段，也以具体问题牵引我逐项改进。课题组为本课题配置了难得的实验条件：高性能计算资源、Franka 协作机械臂及知行末端夹爪、深度相机等感知器件，使许多设想得以工程化验证。正是在控制器接口打通、LeRobot 数据规范对齐以及模型训练与实机部署等环节的支持下，我才把阶段性构想落实为可运行、可评测的系统，顺利完成毕业设计。

  其次要感谢父母与家人，在生活琐事与情绪波动时给予的照料与理解；正是这份稳定依靠，让我衣食无忧，也能更安心地求学与科研。

  我还要感谢同窗与合作者。陈禹彤是本课题最密切的搭档：实验方案论证、示教采集与数据抽检、训练配置对齐以及实机闭环对比与结果整理，我们在一次次并肩推进中学会分工与担当；联调排错中的并肩作战，也成为本科阶段难忘的记忆。刘恒华承担主动臂结构设计及 RealSense 安装件建模打印，关乎装配与视野的细节直接影响平台可用性。感谢机器人队伙伴们，与你们在调试与讨论中的相处让我保持工程敏感与合作信任。亦感谢李思捷、冷龙康、王晨宇、朱培炎等好友，学业建议与生活陪伴让大学四年充实温暖。

  最后感谢母校北京理工大学提供的课堂、实验室与社团土壤。四年里我积累了专业知识与工程方法，也更懂得在压力下自我调节、在分歧中寻求共识。这段经历未必尽善尽美，却足够真诚；它将鞭策我以更踏实的心态走向下一阶段的人生旅途。

  %\textcolor{blue}{致谢正文样式与文章正文相同：宋体、小四；行距：22 磅；间距段前段后均为 0 行。阅后删除此段。}

  % 既然写到这里，那估计就要大功告成了。无论过程与结果如何，总是人生一段独特体验。祝贺！
  % 如有兴趣，可填写以下反馈问卷，帮助 BIThesis 未来优化。
  % https://wj.qq.com/s2/22532808/254f/
\end{acknowledgements}
